\documentstyle[% 


righttag% 


,11pt% 


,amssymb% 


,russian%               remove in Eng. 


,izvru%                 Set izven% in Eng. 


]{amsart} 


 


%       Math definitions 


 


\newcommand{\tts}[1]{} 


 


%\hoffset=-15mm%                  For Eng. 


 


\setcounter{page}{63} 


\god{1999} 


\nomer{10~(449)} %                        Remove in Eng. 


%\nomer{Vol.\,43, No.\,10}%               Set in Eng. 


%\str{\pageref{begin}--\pageref{end}}%  Set in Eng. 


\udk{517.958} 


 


\author{\it Л.Д.\,Эскин} 


% NB:   ^^ remove in Eng. 


\title{Об интегральном уравнении, описывающем ориентационные  


фазовые переходы в модели Парсонса} 


\thanks{Работа выполнена при финансовой поддержке Российского фонда 


фундаментальных исследований (проект 97-01-00375).} 


 


\date{} 


% \pagestyle{myheadings}%         Set in Eng. 


\pagestyle{plain} %              Remove in Eng. 


\begin{document} 


% \thispagestyle{plain}%          Set in Eng. 


\maketitle 


\label{begin} 


 


{\bf 1.} В 1949 г. Л.\,Онзагером (изложение результатов 


Онзагера и их дальнейшее развитие см. в [1]) были подробно исследованы 


термодинамические свойства системы сильно вытянутых жестких 


неполярных цилиндрических стержней ($\delta=dl^{-1}\ll 1$, $d$ --- 


диаметр, 


$l$ --- длина стержня) с парным взаимодействием типа стерического 


отталкивания (модель Онзагера исключенного объема). Было  


показано, что в системе, ориентационно-разупорядоченной при низких 


концентрациях, с увеличением концентрации происходит фазовый переход 


первого рода в анизотропную (ориентационно-упорядоченную)  


фазу, трактуемую как жидкокристаллический нематик. Все 


термодинамические 


свойства изотропной фазы описываются равномерной функцией 


распределения ориентации осей частиц с плотностью $f(n)=1$ ($n$ --- орт 


оси стержня), термодинамические свойства нематика описываются отличной 


от единицы плотностью $f$, имеющей единственный максимум в  


направлении директора (направление преимущественной ориентации 


осей), инвариантной относительно поворотов вокруг этого направления 


и замены $n\to -n$. 


 


Для $f(n)$ из условия минимума свободной энергии системы стержней, 


найденной в приближении второго вириального коэффициента,  


Онзагер получил нелинейное интегральное уравнение, которое 


исследовалось 


(в основном, численными методами) во многих физических  


работах. Достаточно полный обзор полученных в этом направлении  


результатов приводится в обзоре [2] и монографии [3]. Обобщению  


модели Онзагера на случай системы магнитных стержней 


посвящены работы [4], [5].  


 


Необходимо отметить, что Онзагер вычислил свободную энергию  


системы стержней лишь в приближении второго вириального коэффициента, 


т.\,е. в предположении малости концентрации системы. Чтобы  


ориентационный фазовый переход в системе стержней мог произойти  


уже при низкой концентрации, необходимо условие $\delta\ll 1$ --- 


условие  


применимости модели Онзагера. 


 


Модель, свободная от этого ограничения и пригодная для любых  


осесимметричных частиц, была предложена Парсонсом [6]. В наиболее 


интересном случае частиц, имеющих форму эллипсоида вращения с полуосями 


$b\ge a$, интегральное уравнение, полученное Парсонсом для 


ориентационной функции $f(n)$, имеет вид 


\begin{equation} 


\nu+\ln f(n')+\lambda\int B(n,n')f(n)dn=0,\tag{1.1} 


\end{equation} 


где ядро $B(n,n')=(1-\chi^2(nn')^2)^{1/2}$, $nn'$ --- скалярное 


произведение  


ортов $n$ и $n'$, $\chi=(b^2-a^2)(b^2+a^2)^{-1}$, неизвестная  


константа $\nu$ определяется условием нормировки 


\begin{equation} 


\int f(n)dn=1.\tag{1.2} 


\end{equation} 


Интегрирование в (1.1) и (1.2) производится по поверхности сферы,


в сферической системе координат с полярной осью, направленной 


вдоль директора нематика, $dn=(4\pi)^{-


1}\sin\theta\,d\varphi\,d\theta$, $\varphi$, $\theta$ --- 


сферические координаты орта $n$. Параметр $\lambda$ в уравнении 


(1.1) определяется соотношением $\lambda=8(1-\chi^2)^{-1/2}J$, 


где функция $J$ при любом $\chi$ ($0<\chi<1$) является монотонно 


возрастающей 


функцией безразмерной переменной $\eta=\frac{4}{3}\pi a^2bc$ ($c=N/V$ 


--- плотность системы эллипсоидов, $\eta$ --- объемная концентрация 


системы).  


Функция $J$ определяется конкретным выбором модели парного потенциала 


и парной корреляционной функции, относительно которых в модели 


Парсонса предполагается скейлинговый характер их зависимости  


от трансляционных и угловых переменных, что и позволило разделить  


трансляционные и ориентационные степени свободы и получить уравнение 


(1.1). Отметим, что уравнение Онзагера для $f(n)$ является частным 


случаем уравнения (1.1), оно получается, если в уравнении 


(1.1)\; $\lambda=2cdl^2$, а ядро $B=(1-(nn')^2)^{1/2}$. 


 


Поскольку ориентационная функция $f(n)$ описывает нематик, то  


она должна быть решением нелинейного интегрального уравнения  


(1.1), удовлетворяющим, кроме условия нормировки (1.2), еще и  


следующим дополнительным условиям:


 


a) $f(n)$ не зависит от угла $\varphi$ ($f(n)=f(\theta)$), 


 


b) $f(\theta)=f(\pi-\theta)$, следовательно, $f(n)$ разлагается в ряд 


Фурье  


по полиномам Лежандра с четным индексом $P_{2s}(n)$, 


 


c) $f(0)=f(\pi)=\max$, других максимумов $f(n)$ не имеет. 


 


\begin{note} Среди полиномов $P_{2s}$ условию c) удовлетворяет 


только полином $P_2$. 


\end{note} 


 


Это замечание весьма важно для дальнейшего. 


 


В данной работе изучаются близкие к изотропному анизотропные 


решения уравнения (1.1), удовлетворяющие условию нормировки  


(1.2) и условиям a)--c). Для этих решений доказывается с помощью  


методов теории ветвления решений нелинейных уравнений (теории 


Ляпунова--Шмидта [7]), что они в окрестности точки бифуркации 


$\lambda=\lambda_b$ 


разлагаются в сходящийся степенной ряд по целым степеням разности 


$\lambda-\lambda_b$. Этот результат получаем на основе исследования 


диаграммы Ньютона уравнения разветвления для уравнения (1.1). 


 


В заключительной части работы строится эффективный алгоритм  


для указанных разложений. 


\medskip 


 


{\bf 2.} Поскольку нас будут интересовать лишь решения уравнения 


(1.1), описывающие нематик (т.\,е. удовлетворяющие условию 


(1.2) и условиям a), b), c)), то будем рассматривать это уравнение  


в банаховом пространстве $C$ непрерывных функций на сфере  


($\|f\|=\sup|f(n)|$), инвариантных относительно поворотов вокруг  


полярной оси сферической системы координат (т.\,е. зависящих  


лишь от угла $\theta$ между вектором $n$ и полярной осью (условие 


a))\,) и  


при замене $n\to -n$ ($\theta\to\pi-\theta$, условие b)). Учитывая, что 


ядро  


$B(n,n')$ зависит лишь от угла $\alpha$ между векторами $n$ и $n'$ 


(следовательно, 


оно инвариантно при их одновременном повороте), причем 


$B(n,n')=B(-n,n')=B(n,-n')$, а мера $dn$ инвариантна относительно  


вращений, нетрудно доказать, что интегральный оператор 


$$ 


h\to A_\lambda h=\lambda\int B(n,n')h(n)dn 


$$ 


отображает пространство $C$ в себя. 


 


Полагая $f=1+h(n)$, где $|h(n)|$ мал (напомним, что рассматриваются 


лишь ориентационные функции $f$, близкие к изотропным), 


получим из (1.1) нелинейное интегральное уравнение 


\begin{equation} 


\nu+\sum^\infty_{l=1}(-1)^{l-1}l^{-1}h^l+ 


A_\lambda h=0\tag{2.1} 


\end{equation} 


для неизвестной функции $h\in C$, удовлетворяющей условию нормировки 


\begin{equation} 


\int h(n)dn=0.\tag{2.2} 


\end{equation} 


Так как ядро $B$ зависит лишь от 


угла $\alpha$, то $\int B(n,n')dn$ 


не зависит от $n'$, откуда, меняя порядок интегрирования в двукратном  


интеграле $I=\int A_\lambda h\,dn'$, 


найдем в силу условия нормировки (2.2), что $I=0$. Интегрируя 


теперь уравнение (2.1) по $n'$, получим уравнение для $h$ 


\begin{equation} 


h+A_\lambda h+\sum^\infty_{l=2}(-1)^{l-1}l^{-1} 


\Big(h^l-\int h^l dn\Big)=0.\tag{2.3} 


\end{equation} 


Левая часть уравнения (2.3) является интегро-степенным рядом по 


$h$ и $\lambda$, регулярно сходящимся при $\|h\|\le q<1$, 


$|\lambda|\le\lambda_0$ ($\lambda_0>0$


произвольно). Малые решения таких уравнений исследуются в теории 


ветвления решений нелинейных интегральных уравнений --- теории  


Ляпунова--Шмидта [7]. При любом $\lambda$ уравнение (2.3) имеет решение  


$h=0$. В теории Ляпунова--Шмидта доказывается возможность 


существования в окрестности точки бифуркации $\lambda=\lambda_b$ 


ненулевого решения $h_\lambda$, 


стремящегося к нулю при $\lambda\to\lambda_b$. 


 


Ядро $B$ интегрального оператора $A_\lambda$ разлагается в ряд Фурье по  


полиномам Лежандра $P_{2s}(\cos\alpha)$ 


$$ 


B=c_0(\chi)-K_1(n,n'), 


$$ 


где 


\begin{equation} 


\begin{split} 


c_0(\chi)&=\frac{1}{2}\int^1_0(1-\chi^2x^2)^{1/2}dx,\quad 


K_1=\sum^\infty_{k=1}(4k+1)c_k(\chi)P_{2k}(nn'), \\ 


c_k(\chi)&=-\frac{1}{4}\int^1_{-1}(1-\chi^2x^2)^{1/2} 


P_{2k}(x)dx. 


\end{split} 


\tag{2.4} 


\end{equation} 


Согласно известной теореме Гобсона [8] ряд (2.4) для ядра $K_1$ 


сходится 


равномерно по $\alpha$. 


 


Нетрудно убедиться, что любое решение $h\in C$ уравнения (2.3), 


удовлетворяющее условию нормировки (2.2), будет удовлетворять  


и уравнению 


\begin{equation} 


h-\lambda\int K_1(n,n')h(n)dn=\sum^\infty_{l=2} 


(-1)^l l^{-1}\Big(h^l-\int h^ldn\Big).\tag{2.5} 


\end{equation} 


Обратно, любое решение $h\in C$ уравнения (2.5) автоматически 


удовлетворяет 


условию нормировки (2.2), следовательно, и уравнению 


(2.3). Действительно, из (2.4) найдем 


$$ 


\int K_1(n,n')dn'=0. 


$$ 


Интегрируя теперь по $n'$ в обеих частях уравнения (2.5) и меняя 


порядок интегрирования в двукратном интеграле 


$$ 


\int\bigg(\int K_1(n,n')h(n)dn\bigg)dn', 


$$ 


получим для $h$ условие нормировки (2.2). 


 


Для полиномов Лежандра справедлива теорема сложения [8], [9] 


\begin{equation} 


P_l(nn')=P_l(n)P_l(n')+2\sum^l_{m=1} 


\frac{(l-m)!}{(l+m)!}P^m_l(n)P^m_l(n') 


\cos m(\varphi-\varphi')\tag{2.6} 


\end{equation} 


(для сокращения записи снова обозначаем $P^m_l(\cos\theta)=P^m_l(n)$, 


$P^m_l$ --- присоединенные сферические функции, $\varphi$, $\theta$, 


$\varphi'$, $\theta'$ --- сферические 


координаты ортов $n$ и $n'$ соответственно). Поскольку $h\in C$, 


т.\,е. удовлетворяет условиям a) и b), то с учетом соотношения  


(2.6) уравнение (2.5) для $h$ можно переписать в 


виде 


\begin{equation} 


h-\lambda\int K(n,n')h(n)dn=\sum^\infty_{l=2} 


(-1)^l l^{-1}\Big(h^l-\int h^l dn\Big),\tag{2.7} 


\end{equation} 


где ядро 


\begin{equation} 


K(n,n')=\sum^\infty_{k=1}(4k+1)c_k(\chi) 


P_{2k}(n)P_{2k}(n').\tag{2.8} 


\end{equation} 


С помощью разложения (2.8) и соотношения ортогональности для 


полиномов Лежандра найдем 


$$ 


\int K(n,n')P_{2k}(n)dn=c_k(\chi)P_{2k}(n'). 


$$ 


Следовательно, полином $P_{2k}$ является собственной функцией ядра  


$K$, принадлежащей собственному значению $c_k(\chi)$. Из (2.7) 


получаем, что точки бифуркации нелинейного интегрального уравнения 


(2.7) определяются соотношением 


$\lambda^{(k)}_b(\chi)=(c_k(\chi))^{-1}$. 


 


Нетрудно доказать, что при $k\ge 1$ и $0<\chi<1$\; $c_k(\chi)>0$. 


С этой целью воспользуемся формулой Родрига ([9], c.\,1039) для 


полиномов 


Лежандра 


\begin{equation} 


P_m(x)=\frac{1}{2^m m!}\,\frac{d^m(x^2-1)^m}{dx^m}.\tag{2.9} 


\end{equation} 


Подставляя (2.9) в соотношение (2.4) для $c_k$, найдем 


\begin{equation} 


c_k(\chi)=-\frac{1}{2^{2k+2}(2k)!}\int^1_{-1} 


(1-\chi^2x^2)^{1/2}\frac{d^{2k}(x^2-1)^{2k}}{dx^{2k}}dx.\tag{2.10} 


\end{equation} 


Интегрируя в правой части равенства (2.10) по частям $2k$ раз, получим 


\begin{equation} 


c_k(\chi)=-\frac{1}{2^{2k+2}(2k)!}\int^1_{-1} 


(x^2-1)^{2k}\frac{d^{2k}(1-\chi^2x^2)^{1/2}}{dx^{2k}}dx.\tag{2.11} 


\end{equation} 


Очевидно, справедливо неравенство  


\begin{equation} 


\frac{d^{2k}(1-\chi^2x^2)^{1/2}}{dx^{2k}}<0.\tag{2.12} 


\end{equation} 


Из (2.12) и (2.11) и следует указанное неравенство для $c_k(\chi)$. 


 


Далее докажем справедливость неравенств  


\begin{equation} 


c_k(\chi)>c_{k+1}(\chi),\quad 0<\chi<1,\ \; k=1,2,\dotsc\tag{2.13} 


\end{equation} 


Действительно, из соотношения (2.4) для $c_k(\chi)$ получаем с помощью  


известного рекуррентного соотношения для полиномов Лежандра [8], [9] 


$$ 


c_k(\chi)-c_{k+1}(\chi)=\frac{1}{4}\int^1_{-1} 


(1-\chi^2x^2)^{1/2}(P_{2k+2}(x)-P_{2k}(x))dx= 


\frac{4k+3}{8(k+1)(2k+1)}I_k(\chi), 


$$ 


где 


\begin{equation} 


I_k=\int^1_{-1}(1-\chi^2x^2)^{1/2}(x^2-1) 


\frac{d P_{2k+1}}{dx}dx.\tag{2.14} 


\end{equation} 


Снова воспользуемся формулой Родрига (2.9). Подставим (2.9) при 


$m=2k+1$ в правую часть соотношения (2.14), а затем проинтегрируем 


$2k+2$ раза по частям. Получим 


\begin{equation} 


I_k=\frac{1}{2^{2k+1}(2k+1)!}\int^1_{-1}(x^2-1)^{2k+1} 


\frac{d^{2k+2}}{dx^{2k+2}}[(1-\chi^2x^2)^{1/2}(x^2-1)]dx.\tag{2.15} 


\end{equation} 


Разложим выражение в квадратных скобках в последнем интеграле в  


ряд Маклорена. После простых преобразований найдем 


\begin{multline} 


(1-\chi^2x^2)^{1/2}(x^2-1)=-1+\big(1-\tfrac{1}{2}\chi^2\big)x^2- 


\tfrac{1}{2}\chi^2\big(1-\tfrac{1}{4}\chi^2\big)x^4+ \\


+\sum^\infty_{m=3}\frac{(2m-5)!!}{(2m-2)!!}\chi^{2m-2} 


\bigg(\frac{2m-3}{2m}\chi^2-1\bigg)x^{2m}.\tag{2.16} 


\end{multline} 


Поскольку 


$$ 


\frac{2m-3}{2m}\chi^2-1<0\quad (0<\chi<1,\ \; m\ge 3), 


$$ 


то из (2.16) легко следует неравенство  


$$ 


\frac{d^{2k+2}}{dx^{2k+2}}[(1-\chi^2x^2)^{1/2}(x^2-1)]<0, 


\quad k=1,2,\dotsc, 


$$ 


после чего из (2.15) получаем сперва $I_k>0$, а затем и требуемое 


неравенство (2.13) для коэффициентов $c_k$. 


 


Из неравенства (2.13) следует, что каждая из точек бифуркации 


$\lambda^{(k)}_b(\chi)$ при любом $\chi\in(0,1)$ является точкой 


бифуркации с одномерным 


ветвлением. Ниже будем рассматривать лишь решение $h\in C$ 


нелинейного интегрального уравнения (2.7), ответвляющееся от 


изотропного 


решения $h=0$ в точке бифуркации 


$\lambda_b(\chi)=\lambda^{(1)}_b(\chi)$, т.\,к. 


только это решение будет давать ориентационную функцию $f=1+h(n)$, 


удовлетворяющую как условию нормировки (1.2), так и всем трем условиям 


a), b), c). Из дальнейшего будет ясно, что решения, ответвляющиеся 


в точках бифуркации $\lambda^{(k)}_b(\chi)$, $k\ge 2$, не будут 


удовлетворять условию c), 


следовательно, и не будут описывать нематик.  


 


Обозначим $\varphi_k(n)=(4k+1)^{1/2}P_{2k}$, так что  


$$ 


\int\varphi_m(n)\varphi_k(n)dn=\delta_{km} 


$$ 


($\delta_{km}$ --- символ Кронекера), и положим $\lambda=\lambda_b+\mu$, 


$$ 


E(n,n')=\lambda_b K(n,n')-\varphi_1(n)\varphi_1(n') 


$$ 


и 


\begin{equation} 


\xi=\int h(n)\varphi_1(n)dn\tag{2.17} 


\end{equation} 


(в физике величина $\xi$ называется параметром порядка и представляет  


большой самостоятельный интерес).


 


Уравнение (2.7) в новых обозначениях можно переписать в виде 


\begin{equation} 


h-\int E(n,n')h(n)dn=\xi\varphi_1+\mu\int K(n,n') 


h(n)dn+\sum^\infty_{l=2}(-1)^l l^{-1} 


\Big(h^l-\int h^l dn\Big).\tag{2.18} 


\end{equation} 


Поскольку, как мы показали выше, каждая из точек бифуркации является 


точкой одномерного ветвления для уравнения (2.7), то в силу 


леммы Шмидта [7] единица не является собственным значением ядра 


$E(n,n')$, и из теории Ляпунова--Шмидта в силу регулярной сходимости 


интегро-степенного ряда в правой части уравнения (2.18)  


получаем, что уравнение (2.7) при достаточно малых $|\xi|$ и $|\mu|$ 


имеет 


единственное малое решение $h(n)\in C$ (следовательно, $h$ 


удовлетворяет 


условиям a) и b); очевидно, оно удовлетворяет и условию нормировки 


(2.2)), которое представляется в виде равномерно  


сходящегося ряда 


\begin{equation} 


h(n)=\xi\varphi_1+\sum_{r+s=2}\xi^r\mu^s a_{rs}(n),\quad 


\int a_{rs}(n)dn=0.\tag{2.19} 


\end{equation} 


Здесь пока неизвестны и параметр порядка $\xi$ (он выражается через 


искомое решение $h$ с помощью соотношения (2.17)) и коэффициенты 


$a_{rs}(n)$. Подставляя (2.19) в правую часть соотношения (2.17), 


получим 


для $\xi$ уравнение разветвления  


\begin{equation} 


\sum^\infty_{m=2}\cal L_{m0}\xi^m+ 


\sum^\infty_{m=0}\xi^m\sum^\infty_{r=1}\cal L_{mr}\mu^r=0,\tag{2.20} 


\end{equation} 


где 


$$ 


\cal L_{ij}=\int a_{ij}(n)\varphi_1(n)dn 


$$ 


($\xi$ является малым решением уравнения разветвления). 


 


Уравнение (2.20) не содержит слагаемого с первой степенью $\xi$ 


(и нулевой степенью $\mu$), поэтому его решения разлагаются в 


сходящиеся 


при малых $\mu$ ($|\mu|\le\mu_0$) степенные ряды по положительным,  


вообще говоря, дробным, степеням $\mu$ (в теории метода Ньютона [7]  


доказывается, что в этих разложениях показатели степени --- 


рациональные дроби, имеющие конечный общий знаменатель). 


 


Покажем, что ненулевое малое решение $\xi=\xi(\mu)$ уравнения 


разветвления (2.20) 


для интегрального уравнения (2.18) единственно и 


разлагается в сходящийся степенной ряд по целым степеням 


$\mu=\lambda-\lambda_b$ в некоторой окрестности точки бифуркации 


$\lambda_b(\chi)$. С учетом (2.19) 


отсюда будет следовать, что существует единственное ненулевое малое 


решение $h(n)\in C$ уравнения (2.18), разлагающееся в некоторой 


окрестности точки бифуркации $\lambda_b$ по целым положительным 


степеням 


$\lambda-\lambda_b$ с коэффициентами, зависящими лишь от угла $\theta$ 


сферической 


системы координат. Эффективный метод построения этого решения 


укажем в п.\,3. 


 


Нам понадобятся некоторые из коэффициентов $\cal L_{ij}$ уравнения 


разветвления 


(2.20). Прежде всего покажем, что $\cal L_{0j}=0$, $j=1,2,\dotsc$ С 


этой целью подставим решение (2.19) в уравнение (2.18). Положим  


$\xi=0$ в полученном тождестве и сравним коэффициенты при одинаковых 


степенях $\mu$ в левой и правой частях полученных разложений  


по целым степеням $\mu$. Если при этом $a_{02}\ne 0$, то левая часть 


будет  


содержать единственное слагаемое второй степени по $\mu$, а именно 


слагаемое 


$$ 


\Big(a_{02}(n')-\int E(n,n')a_{02}(n)dn\Big)\mu^2. 


$$ 


Поскольку единица не является собственным значением ядра $E$, то 


коэффициент при $\mu^2$ в этом слагаемом будет отличен от нуля. В то 


же время нетрудно заметить, что правая часть в случае $a_{02}\ne 0$ 


будет  


содержать лишь слагаемые не ниже третьей степени $\mu$. Следовательно, 


$a_{02}=0$, а значит, и $\cal L_{02}=0$. Если теперь $a_{03}\ne 0$, то в  


левой части полученного разложения будет содержаться слагаемое 


$$ 


\Big(a_{03}(n')-\int E(n,n')a_{03}(n)dn\Big)\mu^3, 


$$ 


т.\,к. снова коэффициент при $\mu^3$ в этом слагаемом будет отличен 


от нуля. В то же время правая часть будет содержать лишь слагаемые 


не ниже четвертой степени относительно $\mu$. Следовательно, 


$a_{03}=0$, а значит, и $\cal L_{03}=0$. Продолжая эти рассуждения, 


получим $a_{0j}=0$, 


а значит, и $\cal L_{0j}=0$. 


 


Вычислим теперь коэффициент $\cal L_{20}$ уравнения разветвления 


(2.20). С этой целью снова подставим ряд (2.19) в уравнение 


(2.18) и сравним коэффициенты при $\xi^2$ в обеих частях полученного 


тождества. Получим линейное неоднородное интегральное уравнение  


для определения коэффициента $a_{20}$ 


\begin{equation} 


a_{20}(n')-\int E(n,n')a_{20}(n)dn=\frac{1}{2} 


(\varphi^2_1(n')-1).\tag{2.21} 


\end{equation} 


 


Ниже используем формулу Клебша--Гордана ([10], c.\,192) 


\begin{gather} 


P_m P_n=\sum^{m+n}_{l=|m-n|} 


\frac{(2l+1)(l+m-n)!(l-m+n)!(m+n-l)!g!^2} 


{(1+m+n+l)![(g-m)!(g-n)!(g-l)!]^2}P_l,\tag{2.22}\\ 


g=2^{-1}(m+n+l)\notag 


\end{gather} 


(суммирование в (2.22) распространяется лишь на целые $l$ той же 


четности, что и $m+n$). С помощью соотношения (2.22) найдем 


\begin{equation} 


2^{-1}(\varphi^2_1-1)=(\sqrt{5}\varphi_1+\varphi_2)/7.\tag{2.23} 


\end{equation} 


Полиномы Лежандра $P_{2k}$, а следовательно, и $\varphi_k$ являются 


собственными 


функциями ядра $E$, т.\,к. 


\begin{equation} 


\int E(n,n')\varphi_k(n)dn=\lambda_b c_k(\chi) 


\varphi_k(n'),\quad k\ge 2;\quad 


\int E(n,n')\varphi_1(n)dn=0.\tag{2.24} 


\end{equation} 


Поэтому из (2.23) следует, что единственное решение уравнения (2.21) 


(напомним, что  


единица не является собственным значением ядра $E$, и приведенное  


однородное уравнение для (2.21) имеет в  


пространстве $C$ лишь нулевое решение) следует искать в виде 


\begin{equation} 


a_{20}=\alpha_1\varphi_1+\alpha_2\varphi_2.\tag{2.25} 


\end{equation} 


Подставляя (2.25) в уравнение (2.21), с учетом соотношений (2.23) и  


(2.24) найдем 


$$ 


\alpha_1=\sqrt{5}/7,\quad \alpha_2=(7(1-\lambda_b c_2))^{-1}. 


$$ 


Таким образом, оба коэффициента $\alpha_1$, $\alpha_2$ определены и 


положительны  


(в силу неравенств (2.13) и определения точки бифуркации $\lambda_b$ 


справедлива оценка $1-\lambda_b c_2>0$). Теперь окончательно получаем 


$\cal L_{20}=\sqrt{5}/7$. 


 


Чтобы найти $a_{11}(n)$, надо после подстановки ряда (2.19) в уравнение 


(2.18) сравнить в обеих частях полученного тождества коэффициенты 


при произведении $\xi\mu$. В результате получим уравнение для  


определения коэффициента $a_{11}(n)$ 


$$ 


a_{11}-\int E(n,n')a_{11}(n)dn=c_1(\chi)\varphi_1(n')= 


\lambda^{-1}_b\varphi_1(n'), 


$$ 


откуда $a_{11}(n)=\varphi_1(n)/\lambda_b$, 


$\cal L_{11}=\lambda^{-1}_b$. Итак, $\cal L_{0j}=0$, $j\ge 2$, но $\cal 


L_{20}\ne 0$, 


$\cal L_{11}\ne0$ ($\chi\in(0,1)$). 


Отсюда следует,  


что убывающая часть диаграммы Ньютона уравнения разветвления  


(2.20), определяющая его малые решения $\xi=\xi(\mu)$, состоит из 


одного 


отрезка, соединяющего точки $(1,1)$ и $(2,0)$. В этом случае,  


как известно [7], уравнение разветвления (2.20), кроме решения  


$\xi=0$, имеет единственное ненулевое решение $\xi(\mu)$, разлагающееся 


в некоторой окрестности точки $\mu=0$ в сходящийся степенной ряд по 


целым положительным степеням $\mu$ 


\begin{equation} 


\xi=\tau_1\mu+\tau_2\mu^2+\dotsb,\tag{2.26} 


\end{equation} 


где 


$$ 


\tau_1=-\cal L_{11}/\cal L_{20}=-\frac{7}{\sqrt{5}\lambda_b}<0. 


$$ 


После замены параметра $\xi$ в ряде (2.19) решением (2.26) 


получим близкое к изотропному решение $f=1+h(n)$ уравнения Парсонса 


(1.1) в виде сходящегося в некоторой окрестности точки бифуркации 


$\lambda_b$ 


ряда по целым неотрицательным степеням $\mu=\lambda-\lambda_b$. Это  


решение ответвляется от изотропного в точке $\lambda=\lambda_b$ (т.\,е. 


$\lim f(n)=1$, $\lambda\to\lambda_b$), удовлетворяет условию нормировки 


(1.2) и 


условиям a) и b). Из (2.19), (2.26) и замечания 1 следует, поскольку 


$\tau_1<0$, что это решение при достаточно малых $|\mu|$ 


будет удовлетворять условию c) тогда и только тогда, когда $\mu<0$, 


т.\,е. при $\lambda<\lambda_b$. Это означает, что направление 


бифуркации в точке $\lambda_b$ 


левое. Из физической литераторы хорошо известно, что точке бифуркации 


с левым направлением бифуркации соответствует фазовый переход 


первого рода со скачком концентрации. 


 


\begin{note} Аналогично строятся и малые решения $h(n)$ (а вместе 


с ними и $f(n)$), ответвляющиеся от тривиального $h=0$ в точках  


бифуркации $\lambda^{(k)}_b$, $k\ge 2$. Эти решения также будут 


представляться 


в виде рядов (2.19), в которых, однако, в первом слагаемом 


множитель $\varphi_1(n)$ придется заменить на 


$\varphi_k(n)=\sqrt{4k+1}P_{2k}(n)$, соответственно 


изменятся и остальные коэффициенты $a_{rs}(n)$. Так как при $k\ge 2$ 


полином $P_{2k}$ в отличие от $P_2$ имеет локальные максимумы внутри 


отрезка $0\le\theta\le\pi$, то для решений, ответвляющихся от 


тривиального при 


$\lambda=\lambda^{(k)}_b$, $k\ge 2$, условие c) не будет выполняться как 


при $\lambda<\lambda^{(k)}_b$, 


так и при $\lambda>\lambda^{(k)}_b$, следовательно, такие решения не 


будут описывать нематик. 


\end{note}


 


{\bf 3.} Здесь будет указан эффективный алгоритм для построения  


разложения анизотропной плотности $f(n)=1+h(n)$, близкой к изотропной 


и удовлетворяющей условию нормировки (1.2) и условиям a), 


b), c), в ряд по целым степеням $\mu=\lambda-\lambda_b$. Сходимость 


этого разложения 


уже была доказана в п.\,2. С этой целью воспользуемся 


уравнением (2.7), предварительно переписав его в виде 


\begin{equation} 


A h=h-\lambda_b\int K(n,n')h(n)dn= 


\sum^\infty_{l=2}(-1)^l l^{-1}\Big(h^l-\int h^l dn\Big) 


+\mu\int K(n,n')h(n)dn,\tag{3.1} 


\end{equation} 


и положим 


\begin{equation} 


h=\sum^\infty_{i=1}h_i(n)\mu^i.\tag{3.2} 


\end{equation} 


Для определения неизвестных коэффициентов формального степенного 


ряда (3.2) воспользуемся методом неопределенных коэффициентов. 


Подставляя (3.2) в (3.1) и сравнивая коэффициенты при одинаковых 


степенях $\mu$ в обеих частях полученного равенства, найдем для 


определения $h_i$ линейное неоднородное интегральное уравнение 


с симметрическим ядром и правой частью, зависящей лишь от 


коэффициентов $h_1,\dotsc, h_{i-1}$. Условием разрешимости этого 


уравнения 


будет условие ортогональности правой части уравнения к решению 


однородного уравнения $A h=0$. Сравнивая коэффициенты при первой 


степени $\mu$, получим однородное уравнение $A h_1=0$, откуда 


$h_1=\tau_1\varphi_1$. 


Величина $\tau_1$ будет определена при построении следующего 


приближения 


$h_2$, уравнение для которого получим, сравнивая коэффициенты при 


$\mu^2$. Будем иметь с учетом (2.22) уравнение 


\begin{equation} 


A h_2=\tau_1\bigg(c_1+\frac{\sqrt{5}}{7}\tau_1\bigg)\varphi_1+ 


\tau^2_1\varphi_2/7.\tag{3.3} 


\end{equation} 


Коэффициент при $\varphi_1$ в правой части уравнения (3.3) должен быть  


равен нулю (правая часть должна быть ортогональна $\varphi_1$). Отсюда 


$\tau_1=-7c_1/\sqrt{5}$ (если $\tau_1=0$, то и $h=0$). Решение $h_2$ 


уравнения 


следует искать в виде 


\begin{equation} 


h_2=\tau_2\varphi_1+\alpha_{21}\varphi_2.\tag{3.4} 


\end{equation} 


Подставляя (3.4) в уравнение (3.3), находим  


$$ 


\alpha_{21}=\frac{\tau^2_1}{7(1-\lambda_b c_2)} 


$$ 


(напомним, что $1-\lambda_b c_2>0$ ($0<\chi<1$)), а $\tau_2$ найдем при 


построении 


третьего приближения $h_3$, для которого получаем уравнение 


\begin{equation} 


A h_3=\int K(n,n')h_2(n)dn+h_1h_2-\int h_1h_2dn+\frac{1}{3} 


\Big(\int h^3_1dn-h^3_1\Big).\tag{3.5} 


\end{equation} 


С помощью формулы Клебша--Гордана (2.22) найдем 


\begin{equation} 


\begin{split} 


\varphi_1\varphi_2&=\frac{6}{7}\varphi_1+ 


\frac{20\sqrt{5}}{77}\varphi_2+\frac{15}{11}\sqrt{5/13} 


\varphi_3, \\ 


\varphi^3_1&=\frac{\sqrt{5}}{7}\bigg(2+3\sqrt{5}\varphi_1+ 


\frac{36}{11}\varphi_2+\frac{90}{11}\sqrt{1/13} 


\varphi_3\bigg). 


\end{split} 


\tag{3.6} 


\end{equation} 


Соотношения (3.6) позволяют представить правую часть уравнения  


(3.5) в виде линейной комбинации 


$$ 


\beta_{31}(\tau_1,\tau_2)\varphi_1+ 


\beta_{32}\varphi_2+\beta_{33}\varphi_3, 


$$ 


где 


$$ 


\beta_{32}=(22\tau_1\tau_2+20\sqrt{5}\tau_1\alpha_{21}+ 


77c_2\alpha_{21}-12\sqrt{5}\tau^3_1)/77,\quad 


\beta_{33}=\frac{15}{11}\sqrt{5/13}\tau_1 


\Big(\alpha_{21}-\frac{2}{7}\tau^2_1\Big). 


$$ 


Коэффициент $\beta_{31}(\tau_1,\tau_2)$ в силу условия существования 


решения уравнения 


(3.5) следует приравнять нулю, откуда получаем 


$$ 


\tau_2=(5\tau^2_1-6\alpha_{21})/\sqrt{5} 


$$ 


и тем самым коэффициент $h_2$ разложения (3.2) полностью определен.  


 


Решение $h_3$ уравнения (3.5) теперь следует искать в виде 


\begin{equation} 


h_3=\tau_3\varphi_1+\alpha_{32}\varphi_2+\alpha_{33}\varphi_3. 


\tag{3.7} 


\end{equation} 


Коэффициенты $\alpha_{32}$, $\alpha_{33}$ определяются немедленно в 


результате подстановки 


(3.7) в уравнение (3.5). Найдем 


$$ 


\alpha_{32}=\frac{\beta_{32}}{1-\lambda_b c_2},\quad 


\alpha_{33}=\frac{\beta_{33}}{1-\lambda_b c_3} 


$$ 


(здесь снова $1-\lambda_b c_3>0$, $0<\chi<1$). Коэффициент $\tau_3$ 


находится из 


условия существования приближения $h_4$, для которого получаем 


интегральное 


уравнение 


\begin{equation} 


A h_4=s(n)-\int s(n)dn+\int K(n,n')h_3(n)dn,\tag{3.8} 


\end{equation} 


где обозначено 


$$ 


s(n)=4^{-1}h^4_1+2^{-1}h^2_2+h_1h_3-h^2_1h_2. 


$$ 


С помощью формулы Клебша--Гордана и найденных выше выражений для 


$h_1$, $h_2$, $h_3$ убеждаемся, что правая часть уравнения (3.8) 


является линейной комбинацией $\varphi_1$, $\varphi_2$, $\varphi_3$, 


$\varphi_4$. Уравнение для 


определения коэффициента $\tau_3$ получим, приравняв нулю коэффициент 


при $\varphi_1$ в этой линейной комбинации. В результате найдем 


$$ 


\tau_3=\frac{1}{7c_1}\bigg[\sqrt{5}\bigg(\frac{25}{11}\tau^4_1+ 


\tau^2_2+10\alpha^2_{21}\bigg)+6(\alpha_{32}+\tau_2\alpha_{21})- 


\tau^2_1\bigg(15\tau_2+\frac{52}{7}\alpha_{21}\bigg)\bigg]. 


$$ 


Теперь полностью определено третье приближение $h_3$. Аналогично  


можно провести вычисление последующих приближений, но выкладки  


становятся все более громоздкими и мы на них не останавливаемся. 


 


Таким образом, с помощью развитого алгоритма с точностью 


до слагаемого порядка $O(\mu^4)$ найдена ориентационная функция и 


параметр 


порядка 


$$ 


\xi=\tau_1\mu+\tau_2\mu^2+\tau_3\mu^3+O(\mu^4). 


$$ 


 


\begin{note} Полученные в этом пункте формулы в явном 


виде выражают $f(n)$ с точностью до слагаемого порядка $O(\mu^4)$ через 


коэффициенты Фурье $c_1$, $c_2$, $c_3$ ядра $B$. Эти коэффициенты 


задаются 


с помощью интегралов (2.4) при $k=1,2,3$ и также вычисляются 


в явном виде, но получающиеся в результате громоздкие соотношения 


приводить не будем. 


\end{note} 


 


Благодарю за внимание и полезные обсуждения Э.Ю.\,Лернера.  


 


\begin{thebibliography}{99} 


 


\bibitem{1} 


Kayser R.F., Raveche H.J. {\em Bifurcation in Onsager's model 


of the isotropic nematic transition} // Phys. Rev. A. -- 1978. 


-- V.\,17. -- \No\,6. -- P.\,2067--2072. 


\bibitem{2} 


Семенов А.Н., Хохлов А.Р. {\em Статистическая физика жидкокристаллических 


полимеров} // УФН. -- 1988. -- Т.\,156. -- Вып.\,3. -- 


С.\,427--476. 


\bibitem{3} 


Де Жен П.\,Ж. {\em Физика жидких кристаллов}. -- М.: Мир, 1977. -- 400 с. 


\bibitem{4} 


Корнев К.Г., Эскин Л.Д. {\em Фазовые переходы в суспензии иглообразных 


магнитов} // Изв. АН СССР. Сер. физ. -- 1991. -- 


Т.\,55. -- \No\,6. -- С.\,1050--1054. 


\bibitem{5} 


Эскин Л.Д. {\em Об интегральном уравнении теории фазовых переходов 


в системе магнитных стержней} // 


ТМФ. -- 1996. -- Т.\,109. -- \No\,3. -- С.\,427--440. 


\bibitem{6} 


Parsons J.D. {\em Nematic ordering on a system of rods} // 


Phys. Rev. A. -- 1979. -- V.\,19. -- \No\,3. -- Р.\,1225--1230. 


\bibitem{7} 


Вайнберг М.М., Треногин В.А. {\em Теория ветвления решений нелинейных 


уравнений}. -- М.: Наука, 1969. -- 528 с. 


\bibitem{8} 


Гобсон Е.В. {\em Теория сферических и эллипсоидальных функций}. -- М.: 


Ин. лит., 1952. -- 476 с. 


\bibitem{9} 


Градштейн И.С., Рыжик И.М. {\em Таблицы интегралов, сумм, рядов 


и произведений}. -- М.:\linebreak ГИФМЛ, 1962. -- 1100 c. 


\bibitem{10} 


Виленкин Н.Я. {\em Специальные функции и теория представлений  


групп}. -- М.: Наука, 1965. -- 588 с. 


 


 


\end{thebibliography} 


 


\vskip 2cm 


 


\post{\it Казанский государственный}{\it Поступила} 


\post{\it университет}{10.02.1998} 


 


\label{end} 


\end{document} 


